% GENERATED FILE. Edit canonical lessons, then run npm run generate:books.
% canonical-source-hash: fnv1a64:b7a00f81cebf44f3
% canonical-lessons: ES-C740-nauseas, ES-C740-aliviar, ES-C740-vomitar, ES-C740-desmayarse, ES-C740-alergia

\chapter{Me siento mal}
\label{ch:es-me-siento-mal}
\hlchaptermodality{740}

\begin{chapteropening}
\textbf{By the end of this chapter:} \emph{I can say that I feel sick or faint and what relieves it.}

Complete the last lesson of chapter 740: I can say that I feel sick or faint and what relieves it.
\end{chapteropening}

\section[las náuseas]{las náuseas — nausea, sickness}

\label{lesson:ES-C740-nauseas}



\begin{quote}
Before the new one: say the Spanish for to operate, then the Spanish for surgery.
\end{quote}

\subsection*{You'll want to know: las náuseas}
\textbf{las náuseas} — \textquotedblleft{}nausea, sickness\textquotedblright{}.

\begin{quote}
El barco me da náuseas.
\end{quote}

\begin{quote}
\emph{Boats make me feel sick.}
\end{quote}

Usually plural: \emph{tener náuseas}, \emph{dar náuseas}.

\begin{cousinweb}
\emph{náuseas} is a learned borrowing from Latin \emph{nausea}, from Greek \emph{nausía} \textquotedblleft{}seasickness\textquotedblright{} (EtymDB), built on \emph{naûs} \textquotedblleft{}ship\textquotedblright{}.
\end{cousinweb}

\begin{grammarlens}[title={What you've built}]
One more word for telling, explaining and describing at B1.
\end{grammarlens}

\subsection*{Your turn}
\begin{itemize}
\raggedright
  \item \emph{Say it:} \emph{las náuseas}
  \item \emph{Say it:} \emph{las náuseas}, once more
  \item \emph{Say it:} a sentence of your own with \emph{las náuseas}
  \item \emph{Recall:} say the Spanish for to operate, then the Spanish for surgery, then say \emph{las náuseas} again
\end{itemize}

\begin{tcolorbox}[breakable,colback=teal!4,colframe=teal!35!black,title={Before you move on}]
What does las náuseas mean? (\textquotedblleft{}Nausea, sickness\textquotedblright{}.) Say it once more.
\end{tcolorbox}
\section[aliviar]{aliviar — to relieve, to ease}

\label{lesson:ES-C740-aliviar}



\begin{quote}
Before the new one: say the Spanish for surgery, then the Spanish for nausea.
\end{quote}

\subsection*{You'll want to know: aliviar}
\textbf{aliviar} — \textquotedblleft{}to relieve, to ease\textquotedblright{}.

\begin{quote}
Esta pastilla alivia el dolor.
\end{quote}

\begin{quote}
\emph{This tablet relieves the pain.}
\end{quote}

\emph{El alivio} is the relief.

\begin{cousinweb}
\emph{aliviar} is inherited from Latin \emph{alleviāre} \textquotedblleft{}to lighten\textquotedblright{} (EtymDB), built on \emph{levis} \textquotedblleft{}light\textquotedblright{}.
\end{cousinweb}

\begin{grammarlens}[title={What you've built}]
One more word for telling, explaining and describing at B1.
\end{grammarlens}

\subsection*{Your turn}
\begin{itemize}
\raggedright
  \item \emph{Say it:} \emph{aliviar}
  \item \emph{Say it:} \emph{aliviar}, once more
  \item \emph{Say it:} a sentence of your own with \emph{aliviar}
  \item \emph{Recall:} say the Spanish for surgery, then the Spanish for nausea, then say \emph{aliviar} again
\end{itemize}

\begin{tcolorbox}[breakable,colback=teal!4,colframe=teal!35!black,title={Before you move on}]
What does aliviar mean? (\textquotedblleft{}To relieve, to ease\textquotedblright{}.) Say it once more.
\end{tcolorbox}
\section[vomitar]{vomitar — to vomit, to be sick}

\label{lesson:ES-C740-vomitar}



\begin{quote}
Before the new one: say the Spanish for nausea, then the Spanish for to relieve.
\end{quote}

\subsection*{You'll want to know: vomitar}
\textbf{vomitar} — \textquotedblleft{}to vomit, to be sick\textquotedblright{}.

\begin{quote}
El niño vomitó en el coche.
\end{quote}

\begin{quote}
\emph{The boy was sick in the car.}
\end{quote}

\emph{Tener ganas de vomitar} — to feel sick.

\begin{cousinweb}
\emph{vomitar} is a learned borrowing from Latin \emph{vomitāre} (EtymDB), from \emph{vomere} \textquotedblleft{}to throw up\textquotedblright{}.
\end{cousinweb}

\begin{grammarlens}[title={What you've built}]
One more word for telling, explaining and describing at B1.
\end{grammarlens}

\subsection*{Your turn}
\begin{itemize}
\raggedright
  \item \emph{Say it:} \emph{vomitar}
  \item \emph{Say it:} \emph{vomitar}, once more
  \item \emph{Say it:} a sentence of your own with \emph{vomitar}
  \item \emph{Recall:} say the Spanish for nausea, then the Spanish for to relieve, then say \emph{vomitar} again
\end{itemize}

\begin{tcolorbox}[breakable,colback=teal!4,colframe=teal!35!black,title={Before you move on}]
What does vomitar mean? (\textquotedblleft{}To vomit, to be sick\textquotedblright{}.) Say it once more.
\end{tcolorbox}
\section[desmayarse]{desmayarse — to faint}

\label{lesson:ES-C740-desmayarse}



\begin{quote}
Before the new one: say the Spanish for to relieve, then the Spanish for to vomit.
\end{quote}

\subsection*{You'll want to know: desmayarse}
\textbf{desmayarse} — \textquotedblleft{}to faint\textquotedblright{}.

\begin{quote}
Se desmayó por el calor.
\end{quote}

\begin{quote}
\emph{She fainted because of the heat.}
\end{quote}

\emph{El desmayo} — the faint.

\begin{cousinweb}
\emph{desmayarse} is the reflexive of \emph{desmayar}, which is usually traced to Old French \emph{esmaier} \textquotedblleft{}to trouble, to dismay\textquotedblright{}; \emph{desmayo} \textquotedblleft{}a faint\textquotedblright{} is built on the verb, and English \emph{dismay} is a cousin.
\end{cousinweb}

\begin{grammarlens}[title={What you've built}]
One more word for telling, explaining and describing at B1.
\end{grammarlens}

\subsection*{Your turn}
\begin{itemize}
\raggedright
  \item \emph{Say it:} \emph{desmayarse}
  \item \emph{Say it:} \emph{desmayarse}, once more
  \item \emph{Say it:} a sentence of your own with \emph{desmayarse}
  \item \emph{Recall:} say the Spanish for to relieve, then the Spanish for to vomit, then say \emph{desmayarse} again
\end{itemize}

\begin{tcolorbox}[breakable,colback=teal!4,colframe=teal!35!black,title={Before you move on}]
What does desmayarse mean? (\textquotedblleft{}To faint\textquotedblright{}.) Say it once more.
\end{tcolorbox}
\section[la alergia]{la alergia — the allergy}

\label{lesson:ES-C740-alergia}



\begin{quote}
Before the new one: say the Spanish for to vomit, then the Spanish for to faint.
\end{quote}

\subsection*{You'll want to know: la alergia}
\textbf{la alergia} — \textquotedblleft{}the allergy\textquotedblright{}.

\begin{quote}
Tengo alergia al polen.
\end{quote}

\begin{quote}
\emph{I am allergic to pollen.}
\end{quote}

\emph{Tener alergia a algo}; the adjective is \emph{alérgico}.

\begin{cousinweb}
\emph{alergia} was borrowed from German \emph{Allergie} (EtymDB), a word an Austrian doctor coined in 1906 from Greek \emph{állos} \textquotedblleft{}other\textquotedblright{} and \emph{érgon} \textquotedblleft{}work\textquotedblright{}.
\end{cousinweb}

\begin{grammarlens}[title={What you've built}]
One more word for telling, explaining and describing at B1.
\end{grammarlens}

\subsection*{Your turn}
\begin{itemize}
\raggedright
  \item \emph{Say it:} \emph{la alergia}
  \item \emph{Say it:} \emph{la alergia}, once more
  \item \emph{Say it:} a sentence of your own with \emph{la alergia}
  \item \emph{Recall:} say the Spanish for to vomit, then the Spanish for to faint, then say \emph{la alergia} again
\end{itemize}

\begin{tcolorbox}[breakable,colback=teal!4,colframe=teal!35!black,title={Before you move on}]
What does la alergia mean? (\textquotedblleft{}The allergy\textquotedblright{}.) Say it once more.
\end{tcolorbox}
